% This template can be used to generate a CV from a Markdown file with a
% YAML frontmatter. This template was created by Hendrik Erz <hendrik@zettlr.com>

% Several instructions in this template come essentially directly from the
% default LaTeX template Pandoc uses. This ensures that it mostly corresponds
% with Pandoc and should be kept updated over time.

% Include some Pandoc specific data that we require in order to support all
% Pandoc features.
$passoptions.latex()$

% We're using the KOMA Article class
\documentclass[11pt,a4paper]{scrartcl}

% Before we do anything else, define the left margin which determines both the
% overall text indent as well as the available space for the label items.
\newlength{\labelmargin}
\setlength{\labelmargin}{3cm}

% Let the user redefine it in the frontmatter
$if(labelmargin)$
\setlength{\labelmargin}{$labelmargin$}
$endif$

% Now, based on the custom label margin, we can define our actual label width:
\newlength{\cvlabelwidth}
\usepackage{calc} % Required package to calculate new lengths
\setlength{\cvlabelwidth}{\labelmargin-0.5cm}

% Also, we need to pull back the left footer so that it is not aligned with the
% main text, but rather has the same 2cm margin from the left edge as the right
% footer has from the right.
\newlength{\cvfooteroffset}
\setlength{\cvfooteroffset}{\labelmargin-2cm} % The 2cm must be the same as right margin

% Margin options (should be defined as soon as possible so that other packages
% have access to the calculations and can calculate their own widths
% accordingly.)
\usepackage[
  left=\labelmargin, % needs to be bigger to account for the outdented item labels
  right=2cm,
  top=2cm,
  bottom=3cm, % Needs to be bigger to have the footer not cramped to the bottom
  nomarginpar
]{geometry}

% Some packages useful for the font handling
\usepackage{fontspec}
\defaultfontfeatures{Mapping=tex-text}
\usepackage{xunicode}
\usepackage{xltxtra}

% sectsty allows us to modify how the headings look.
\usepackage{sectsty}

% The xcolor package provides some color functionality
\usepackage{xcolor}

% Additional packages
\usepackage{amsmath}
\usepackage{amsfonts}
\usepackage{amssymb}
\usepackage{graphicx}

% This is used to use the adjustwidth package which allows us to define a
% different margin for the header of the CV than for the CV sections.
\usepackage{changepage}

% Icons
\usepackage{fontawesome6}
\faStyle{regular}

%%%%%%%%%%%%%%%%%%%%%%%%%%%%%%%%%%%%%%%%%%%%%%%%%%%%%%%%%%%%%%%%%%%%%%%%%%%%%%%%
% Additional Pandoc-specific includes required to support most Pandoc features.
$fonts.latex()$
$font-settings.latex()$
$common.latex()$

%%%%%%%%%%%%%%%%%%%%%%%%%%%%%%%%%%%%%%%%%%%%%%%%%%%%%%%%%%%%%%%%%%%%%%%%%%%%%%%%

% Define default colors
\definecolor{AccentColor}{rgb}{0.5,0.0,0.4}
\definecolor{TextColor}{rgb}{0.2,0.2,0.2}
\definecolor{MutedTextColor}{rgb}{0.5,0.5,0.5}

% Allow the user to redefine the colors
$if(accentColor)$
\definecolor{AccentColor}{rgb}{$accentColor$}
$endif$
$if(textColor)$
\definecolor{TextColor}{rgb}{$textColor$}
$endif$
$if(mutedColor)$
\definecolor{MutedTextColor}{rgb}{$mutedColor$}
$endif$

% The sections being used will all be colored in accent
\partfont{\color{AccentColor}}
\sectionfont{\color{TextColor}}
\subsectionfont{\color{TextColor}}

% Document fonts. Uses TeX default, unless the user overrides it.
$if(mainfont)$
\setmainfont{$mainfont$}
$endif$
$if(sansfont)$
\setsansfont{$sansfont$}
$endif$

% i18n / l10n
\usepackage{polyglossia}
\setdefaultlanguage{english}

% Allow the user to set the language (stolen from the Pandoc template)
$if(lang)$
\setmainlanguage[$for(polyglossia-lang.options)$$polyglossia-lang.options$$sep$,$endfor$]{$polyglossia-lang.name$}
$for(polyglossia-otherlangs)$
\setotherlanguage[$for(polyglossia-otherlangs.options)$$polyglossia-otherlangs.options$$sep$,$endfor$]{$polyglossia-otherlangs.name$}
$endfor$
$endif$

% For better page headers and footers
\usepackage{fancyhdr}
\pagestyle{fancy}
\fancyhead{} % No page header
\fancyfoot[L]{\small\color{MutedTextColor}$if(name)$$name$$endif$} % Left Author
\fancyfoot[C]{\small\color{TextColor}Page \thepage} % Centered page number
\fancyfoot[R]{\small\color{MutedTextColor}Last updated: \today} % Right last updated
% No horizontal rules between head/foot and main body
\renewcommand{\headrulewidth}{0pt}
\renewcommand{\footrulewidth}{0pt}
\fancyfootoffset[L]{\cvfooteroffset}

% Obligatory URL stuff: Make sure links look like normal
% text and can be broken up to fit the page flow.
\usepackage[obeyspaces,spaces]{url}
\usepackage[breaklinks,hidelinks]{hyperref}
% This command makes URLs look like the rest of the text
\urlstyle{same}

% The package enumitem is the central package for the CV sections. It is used to
% display a label in our accent color and also aligns the text properly. We're
% choosing a very narrow label width, so keep labels brief. It is intended only
% for year numbers. I have validated that this works with single years (2026),
% year-ranges (2024-2026), and "Since" labels (Since 2024). Anything longer may
% reach into the printer margins or even get cut off from the page. If you
% really need bigger labels, adjust these numbers.
% Please also note, if you're customising this, that labelsep should
% be the opposite of itemindent, since otherwise all lines except the
% first one of an item's "text" will be wrongly indented.
\usepackage{enumitem}
% Documentation for allowed values: https://mirrors.hust.edu.cn/CTAN/macros/latex/contrib/enumitem/enumitem.pdf
\setlist[description]{
  % Label spacing:
  labelwidth=\cvlabelwidth, % Width of the label description
  labelsep=0.5cm, % Space between label and item
  align=right, % Label alignment
  labelindent=-\labelmargin, % Pull the labels into the margin
  % Item values spacing:
  leftmargin=0cm, % Make the item values align with the headings/page margin
  itemindent=!, % ! means: calculate from other values
  font=\color{AccentColor}
}
% TIP: The package makes the following horizontal spacing assumptions (taken
% from enumitem, p. 6):
% <labelindent> <labelwidth> <labelsep - itemindent>   <itemindent>
% [                    leftmargin                  ]
% RELATIONSHIP BETWEEN VALUES:
% leftmargin + itemindent = labelindent + labelwidth = labelsep

% Remove all heading numberings (even from the custom body-text)
\makeatletter
\renewcommand{\@seccntformat}[1]{}
\makeatother

% Metadata for the PDF
\author{$if(name)$$name$$endif$}
\title{$if(name)$$name$$endif$ - CV}

%%%%%%%%%%%%%%%%%%%%%%%%%%%%%%%%%%%%%%%%%%%%%%%%%%%%%%%%%%%%%%%%%%%%%%%%%%%%%%%%
\begin{document}

% Immediately switch to text color
\color{TextColor}

% Make the header of the CV have the standard 2cm margin. NOTE: This command
% adds and subtracts the numbers from the existing margins, it does not define
% absolute margins.
\begin{adjustwidth}{-1cm}{0cm}

% Main heading is the name. Warn the user if they forgot the property.
\part*{$if(name)$$name$$else$Please provide a name!$endif$}

% We're using "minipages" to display the address versus contact info
% in a two-column layout. Note that the end of the first minipage
% MUST be on the same line as the begin of the second (which we do
% by using a percent-sign IMMEDIATELY after the end.)
\noindent
$if(address)$
  % We have an address section -> two minipages to have two blocks
  \begin{minipage}{0.6\textwidth}
  $if(occupation)$\textbf{$occupation$} \\\\$endif$% No newline here
  $for(address)$$address$ $sep$ \\$endfor$ % Print the address line(s)
  \end{minipage}% NOTE the load-bearing comment
  \begin{minipage}{0.5\textwidth}
  % Print contact info if there is some.
  $if(contact)$
  \begin{tabular}{ r l }
    & \textbf{Contact} \\ [2ex] % Adds some space
    $if(contact.email)$\faEnvelope & \href{mailto:$contact.email$}{$contact.email$}\\ [0.5ex]$endif$
    $if(contact.phone)$\faPhone & \href{tel:$contact.phone$}{$contact.phone$} \\ [0.5ex]$endif$
    $if(contact.orcid)$\faOrcid & \href{https://orcid.org/$contact.orcid$}{$contact.orcid$}\\ [0.5ex]$endif$
    $if(contact.scholar)$\faGoogleScholar & \href{https://scholar.google.com/citations?user=$contact.scholar$}{Scholar} \\ [0.5ex]$endif$
    $if(contact.website)$\faGlobe & $if(contact.website_title)$\href{$contact.website$}{$contact.website_title$}$else$\url{$contact.website$}$endif$\\ [0.5ex]$endif$
    $if(contact.bluesky)$\faBluesky & \href{https://bsky.app/profile/@$contact.bluesky$}{@$contact.bluesky$}\\ [0.5ex]$endif$
    $if(contact.twitter)$\faTwitter & \href{https://www.twitter.com/$contact.twitter$}{@$contact.twitter$}\\ [0.5ex]$endif$
    $if(contact.mastodon)$\faMastodon & $if(contact.mastodon_title)$\href{$contact.mastodon$}{$contact.mastodon_title$}$else$\url{$contact.mastodon$}$endif$\\$endif$
  \end{tabular}
  $endif$
  \end{minipage}
$elseif(contact)$
  % We don't have an address block, only contact -> single-line
  \begin{center}
  \small
  $if(contact.email)$\faEnvelope{} \href{mailto:$contact.email$}{$contact.email$}$endif$%
  $if(contact.phone)$\, \faPhone{} \href{tel:$contact.phone$}{$contact.phone$}$endif$%
  $if(contact.orcid)$\, \faOrcid{} \href{https://orcid.org/$contact.orcid$}{$contact.orcid$}$endif$%
  $if(contact.scholar)$\, \faGoogleScholar{} \href{https://scholar.google.com/citations?user=$contact.scholar$}{Scholar}$endif$%
  $if(contact.website)$\, \faGlobe{} $if(contact.website_title)$\href{$contact.website$}{$contact.website_title$}$else$\url{$contact.website$}$endif$$endif$%
  $if(contact.bluesky)$\, \faBluesky{} \href{https://bsky.app/profile/@$contact.bluesky$}{@$contact.bluesky$}$endif$%
  $if(contact.twitter)$\, \faTwitter{} \href{https://www.twitter.com/$contact.twitter$}{@$contact.twitter$}$endif$%
  $if(contact.mastodon)$\, \faMastodon{} $if(contact.mastodon_title)$\href{$contact.mastodon$}{$contact.mastodon_title$}$else$\url{$contact.mastodon$}$endif$$endif$%
  \end{center}
$endif$

% Here comes the custom body text (the text from the Markdown document)
$body$

% Separate the CV header from the sections with a ruler.
{\color{AccentColor}\rule{\linewidth}{0.2mm}}

% End special margins area here
\end{adjustwidth}

% Now come the sections
$for(sections)$
\section*{$sections.title$}

% NOTE: The "description" environment is being set up in the header with the
% setlist-command.
\begin{description}

$for(sections.items)$
% DEBUG: Use the following command to draw the various spacings for debugging
% purposes: \DrawEnumitemLabel
% The command draws two vertical lines to indicate the leftmargin, and four
% horizontal bars, in this order: labelindent, labelwidth, labelsep, itemindent
\item [$if(sections.items.label)$$sections.items.label$$endif$] $for(sections.items.text)$$sections.items.text$$sep$ \newline $endfor$ 
$endfor$

\end{description}
$endfor$

\end{document}
